\documentclass[conference]{IEEEtran}

\usepackage{cite}
\usepackage{amsmath,amssymb,amsfonts}
\usepackage{algorithm}
\usepackage{algpseudocode}
\usepackage{graphicx}
\usepackage{booktabs}
\usepackage{multirow}
\usepackage{xcolor}
\usepackage{url}

% ---------------------------------------------------------------------------
% Placeholder macros. Set \finaltrue to hide all of them in a submission build.
% ---------------------------------------------------------------------------
\newif\iffinal \finalfalse
\iffinal
  \newcommand{\todo}[1]{}
  \newcommand{\X}{}
\else
  \newcommand{\todo}[1]{\textcolor{red}{[TODO: #1]}}
  \newcommand{\X}{\textcolor{red}{XX}}
\fi
\newcommand{\figph}[2]{%
  \fbox{\parbox[c][#1][c]{0.95\linewidth}{\centering\small\textcolor{red}{#2}}}}

\newcommand{\code}[1]{\texttt{#1}}
\newcommand{\pre}{\prec}
\newcommand{\one}{\mathbf{1}}
\newtheorem{proposition}{Proposition}
\newtheorem{remark}{Remark}

\begin{document}

\title{Interleaving Multiple Spoken Commands\\for Service Robots under Execution Failure}

\author{
\IEEEauthorblockN{\todo{Author 1}, \todo{Author 2}, \todo{Author 3}}
\IEEEauthorblockA{\todo{Affiliation}\\
\todo{email}}
}

\maketitle

\begin{abstract}
Service robots are increasingly asked to accept several instructions at once
and decide for themselves how to carry them out. The RoboCup@Home General
Purpose Service Robot (GPSR) test rewards this explicitly: a robot that
receives three commands together and executes them in an interleaved order
earns a bonus. Interleaving reduces travel, but it also changes what the robot
carries, when, and which information is available when it answers a question.
We present a planner that merges commands parsed by a language model into a
single plan, using a dependency analysis to decide which parts of different
commands may be reordered and an exact dynamic program over per-command
progress to minimize expected execution time under a single-gripper
constraint. When an action fails, the robot replans from its current state
instead of restarting each command. We also introduce GPSR-Sim, a benchmark
built from the official command generator, with \X{} home layouts, stochastic
skill models and success checks on the final world state. Across \X{} command
batches, the planner reduces task time by \X\% relative to one-by-one
execution and by \X\% relative to back-to-back execution, while \X{} answer
errors caused by naive interleaving are eliminated. \todo{Robot result
sentence from RoboCup 2026.}
\end{abstract}

\begin{IEEEkeywords}
task planning, service robots, RoboCup@Home, behaviour trees, language models
\end{IEEEkeywords}

% ===========================================================================
\section{Introduction}
% ===========================================================================

Domestic service robots are usually evaluated one instruction at a time. In
practice, people give several requests together (``bring me a drink, then
tell me how many people are in the kitchen, and put the book back on the
shelf''), and the robot must decide in what order to do them. The
RoboCup@Home GPSR test~\cite{rulebook} formalizes this setting. The robot
receives three commands, chooses whether to hear them one at a time or all at
once, and receives a 200-point bonus if it completes them in an interleaved
order that saves time or motion. Commands are produced by an official
generator~\cite{cmdgen} and rephrased by a language model before being spoken,
so the robot must also parse free-form speech.

Ordering several commands so as to minimize travel resembles a
pickup-and-delivery routing problem~\cite{savelsbergh1995,frederickson1978},
and for three commands the search space is small. The difficulty lies
elsewhere. First, commands are not independent. A command that moves an
object can change the answer to another command that inspects it, and
results stored for one command can be read by another if the executor does not
keep them apart. In our earlier competition system, using sequential commands caused the
robot to report \todo{N} incorrect answers in \todo{M} simulated batches
(Section~\ref{sec:exp}). Second, skills fail. A missed detection in the middle
of an interleaved plan leaves the robot holding an object for another command,
and a recovery strategy that restarts commands in their original order can
violate the same constraints the planner was designed to respect. Third, the
quantity that matters is the score obtained within a time limit, not the
distance driven.

\todo{Thing is that before it was purely sequential (Luis)}

This paper makes the following contributions:
\begin{itemize}
  \item A formulation of multi-command interleaving in which parsed commands
    are segmented, annotated with the state they read and write, and combined
    under a partial order that only permits reordering of non-interfering
    segments (Section~\ref{sec:problem}).
  \item A dynamic program over per-command progress vectors, whose state
    space grows with $\prod_c (L_c+1)$ rather than $2^M$, with a proof that it
    is exact for any objective made of per-step costs, and extensions to the
    interleaving bonus and the time limit (Section~\ref{sec:method}).
  \item A replanning executor that resumes from the robot's current pose and
    gripper state after a failure, and an exact formulation of the problem with
    failures as an acyclic Markov decision process, which we use to measure
    how far replanning is from the best achievable policy
    (Section~\ref{sec:mdp}).
  \item GPSR-Sim, an open benchmark of command batches drawn from the official
    generator, executed in \X{} home layouts with stochastic skill outcomes
    and judged on the final world state (Section~\ref{sec:bench}).
  \item An evaluation against one-by-one execution, back-to-back execution,
    greedy ordering, a classical planner and a language-model planner, plus
    results from the RoboCup 2026 competition (Section~\ref{sec:exp}).
\end{itemize}
\todo{Add a teaser figure: merged vs.\ sequential trajectories with a Gantt bar per command.}

% ===========================================================================
\section{Related Work}
% ===========================================================================

\subsection{Language models for robot task planning}
Language models are now widely used to turn instructions into robot plans,
either by scoring skills against affordances~\cite{saycan}, by generating
programs~\cite{progprompt}, or by translating goals into PDDL for a classical
planner~\cite{llmp}. Behaviour trees have also been generated directly by
language models~\cite{btgenbot}. These systems typically handle one
instruction at a time. \todo{Add work on multi-instruction or multi-goal
LLM planning, if any is found.}

\subsection{Routing with precedence and capacity}
Visiting a set of locations with pickups that must precede deliveries and a
vehicle that carries one item at a time is the stacker crane
problem~\cite{frederickson1978}, a special case of the general
pickup-and-delivery problem~\cite{savelsbergh1995}. Exact dynamic programming
for sequencing problems dates back to Held and Karp~\cite{heldkarp1962}. We
use this literature for the travel component and focus on the parts it does
not model: interference between commands and execution failures.

\subsection{RoboCup@Home and embodied benchmarks}
Recent GPSR systems use foundation models for parsing and
planning~\cite{nimbro2024,gpsrfm2023,prompteng2024}, but to our knowledge none
evaluates interleaved execution of several commands. Embodied planning
benchmarks such as PARTNR~\cite{partnr} and EmbodiedBench~\cite{embodiedbench},
built on simulators such as Habitat~3.0~\cite{habitat3}, evaluate long-horizon
plans but not the competition setting, where commands arrive by speech,
execution is timed, and partial credit is awarded. \todo{Check for RoboCup@Home
simulation league work (e.g.\ SIGVerse / DSPL simulator) and cite.}

% ===========================================================================
\section{Problem Formulation}\label{sec:problem}
% ===========================================================================

\subsection{Commands and segments}
A batch consists of $n$ commands. A language-model parser~\todo{cite or
describe BAML schema} turns command $c$ into a sequence of actions
$a_{c,1},\dots,a_{c,L'_c}$ drawn from a fixed set of $K=13$ skills
(Table~\ref{tab:skills}). We split each sequence into \emph{segments}: a segment
starts with a \code{go\_to} and contains every action up to the next
\code{go\_to}. Command $c$ then has segments $s_{c,1},\dots,s_{c,L_c}$. Each
segment has an entry location $e(s)$, the target of its \code{go\_to}, and an
exit location $x(s)$, which equals $e(s)$ unless the segment contains a
\code{guide\_person\_to} or \code{follow\_person\_until}, in which case it is
that action's destination. A segment with no leading \code{go\_to} is attached
to the preceding segment of the same command, or to the robot's position when
the command was given.

\begin{table}[t]
\centering
\caption{Skills produced by the parser and their state effects.}
\label{tab:skills}
\resizebox{\columnwidth}{!}{%
\begin{tabular}{lll}
\toprule
Skill & Reads & Writes \\
\midrule
\code{go\_to} & -- & robot pose \\
\code{pick\_object} & object pose & gripper, object pose \\
\code{place\_object}, \code{give\_object} & gripper & gripper, object pose \\
\code{find\_person}, \code{find\_person\_by\_name} & person pose & attention \\
\code{get\_person\_info} & attention & register$_c$ \\
\code{count}, \code{get\_visual\_info} & object poses & register$_c$ \\
\code{follow\_person\_until}, \code{guide\_person\_to} & attention & pose, person pose \\
\code{answer\_question} & -- & -- \\
\code{say\_with\_context} & register$_c$ & -- \\
\bottomrule
\end{tabular}}
\todo{Verify read/write sets against the skill implementations.}
\end{table}

\subsection{Interference between commands}
Each action $a$ has a read set $R(a)$ and a write set $W(a)$ over state
variables: the gripper, the person currently attended to, a result register
per command, and the pose of each object and person mentioned in the batch.
For a segment, $R(s)$ and $W(s)$ are the unions over its actions. Two segments
$s$ and $s'$ from different commands \emph{interfere} if
\begin{equation}
W(s)\cap\big(R(s')\cup W(s')\big)\neq\emptyset
\;\;\text{or}\;\;
W(s')\cap R(s)\neq\emptyset .
\end{equation}
Result registers are indexed by command, so information-gathering actions of
different commands never interfere through them. The gripper is handled
separately below. We define the partial order $\pre$ as the transitive closure
of (i) $s_{c,j}\pre s_{c,j+1}$ for every command, and (ii) $s\pre s'$ for every
interfering pair, oriented by the order in which the commands were issued.
\todo{Discuss whether issue order is the right default, or whether the operator
should be asked.}

\subsection{Feasible schedules}
A schedule $\pi$ is a linear extension of $\pre$ over all segments. It is
feasible if it respects the gripper capacity: whenever command $c$ holds an
object, the next segment must belong to $c$ or must not contain a gripper
action. Holding is determined by the last gripper action executed, so a segment
of the form \code{place}, \code{pick} ends holding.

\subsection{Objective}
Let $d(\cdot,\cdot)$ be the path length on the map, $v$ the nominal speed and
$\tau(s)$ the expected duration of the non-navigation actions in $s$. Each
navigation also has a fixed overhead $\kappa$ (stowing the arm, announcing the
destination), which is paid only when the robot actually has to move. The
expected task time is
\begin{equation}
T(\pi)=\sum_{i=1}^{M}\frac{d\big(x(\pi_{i-1}),e(\pi_i)\big)}{v}
+\kappa\,\one\big[e(\pi_i)\neq x(\pi_{i-1})\big]+\tau(\pi_i),
\end{equation}
with $x(\pi_0)=p_0$ the robot's pose when planning starts. Our competition
system removed consecutive \code{go\_to}s to the same place after planning;
including $\kappa$ in the cost lets the planner account for that saving
instead.

Minimizing $T$ ignores risk. We also penalize the time spent carrying an
object, since a failure while holding is the costliest kind. Let $h_{i-1}$
denote the command holding the gripper before step $i$ (or $\varnothing$). The
holding time accumulated in step $i$ is
$\one[h_{i-1}\neq\varnothing]\,d(x(\pi_{i-1}),e(\pi_i))/v+\rho(\pi_i,h_{i-1})$,
where $\rho(s,h)$ is the time spent holding inside segment $s$, which depends
only on $s$ and on whether an object is held when $s$ starts. The objective is
the sum of per-step costs
\begin{equation}\label{eq:stage}
J(\pi)=\sum_{i=1}^{M} g\big(x(\pi_{i-1}),h_{i-1},\pi_i\big),
\end{equation}
\begin{multline*}
g(x,h,s)=\Big(1+\lambda\one[h\neq\varnothing]\Big)\frac{d(x,e(s))}{v}
+\kappa\,\one[e(s)\neq x]\\+\tau(s)+\lambda\rho(s,h),
\end{multline*}
where $\lambda$ trades time against exposure to failure while holding. Writing
$J$ as a sum of costs that depend only on the current pose, the gripper holder
and the next segment is what makes the dynamic program in
Section~\ref{sec:method} exact (Proposition~\ref{prop:exact}).
\todo{Derive $\lambda$ from the score sheet if possible, or select it on a
held-out set of batches.}

% ===========================================================================
\section{Method}\label{sec:method}
% ===========================================================================

\subsection{Dynamic program over progress vectors}
Because segments of one command are totally ordered, the set of segments
executed so far is fully determined by the vector
$\mathbf{k}=(k_1,\dots,k_n)$, where $k_c$ is the number of segments of command
$c$ already executed. A state is a pair $(\mathbf{k},\ell)$, where $\ell$ is
the command that executed the last segment. We add a virtual start state
$(\mathbf{0},0)$ whose pose is $p_0$, and write $x(\mathbf{k},\ell)$ for the
robot's pose in state $(\mathbf{k},\ell)$, i.e.\ $x(s_{\ell,k_\ell})$, or $p_0$
in the start state. The gripper holder $h(\mathbf{k})$ and the set
$E(\mathbf{k})$ of commands whose next segment has all its $\pre$-predecessors
executed and is allowed by the gripper rule are functions of $\mathbf{k}$
alone. The cost-to-go satisfies
\begin{equation}\label{eq:bellman}
V(\mathbf{k},\ell)=\min_{c\in E(\mathbf{k})}
\Big[g\big(x(\mathbf{k},\ell),h(\mathbf{k}),s_{c,k_c+1}\big)
+V(\mathbf{k}+\mathbf{e}_c,c)\Big],
\end{equation}
with $V(\mathbf{L},\cdot)=0$ for $\mathbf{L}=(L_1,\dots,L_n)$. The number of
states is at most $1+n\prod_c(L_c+1)$ and each has at most $n$ successors, so
the running time is $O(n^2\prod_c(L_c+1))$. For a GPSR batch ($n=3$,
$L_c\le4$) this is at most $3\cdot5^3=375$ states. Algorithm~\ref{alg:dp}
gives the procedure.

\begin{proposition}\label{prop:exact}
For any objective of the form~\eqref{eq:stage}, the recursion~\eqref{eq:bellman}
evaluated at $(\mathbf{0},0)$ returns the minimum of $J$ over all feasible
schedules, and following its minimizers yields an optimal schedule.
\end{proposition}

\begin{IEEEproof}[Proof sketch]
Given $(\mathbf{k},\ell)$, the executed segments are exactly
$\{s_{c,j}: j\le k_c\}$, so the remaining segments, the eligible set
$E(\mathbf{k})$ and the gripper holder are determined. The holder is
determined because each command's gripper state after its first $k_c$ segments
is fixed, and the capacity constraint allows at most one holder. The robot's
pose is $x(\mathbf{k},\ell)$. The set of feasible completions from a state and
the cost of each completion therefore depend only on $(\mathbf{k},\ell)$, not on
how the state was reached. This is the principle of optimality, so an optimal
schedule restricted to its suffix from any state is optimal for that state, and
\eqref{eq:bellman} holds. Every transition increases $\sum_c k_c$ by one, so the
state graph is acyclic. Since $\mathbf{k}+\mathbf{e}_c$ is lexicographically
larger than $\mathbf{k}$, processing states in reverse lexicographic order
evaluates every successor before its predecessors, and a single backward pass
computes $V$ exactly.
\end{IEEEproof}

\begin{remark}
Optimality is relative to the model in two ways. Segments are treated as atomic,
so schedules that execute another command's actions in the middle of a segment
are not considered. Same-location segments of different commands are still
grouped when that helps, since moving between them costs nothing. The result
also holds only for the cost model: a cost that the planner does not see, such
as the navigation overhead $\kappa$ when it is omitted, can make the computed
schedule suboptimal on the robot.
\end{remark}

\subsection{Objectives that are not additive}
Two quantities in the GPSR score sheet do not decompose into per-step costs.
Both can be handled by adding a variable to the state, after which
Proposition~\ref{prop:exact} applies to the larger state space.

\emph{Interleaving bonus.} The bonus depends on whether the schedule
interleaved commands at all. We add a bit $b$ that becomes 1 when a segment of
command $c$ is executed while another command $c'$ has
$0<k_{c'}<L_{c'}$, and subtract the bonus $B$ from the terminal cost when
$b=1$. This doubles the number of states. \todo{Check with the technical
committee what counts as ``meaningful'' interleaving.}

\emph{Time limit.} The test ends after $T_{\max}=420$\,s. Maximizing the
probability of finishing, or the expected score when unfinished commands earn
nothing, depends on the total elapsed time rather than a sum of costs. We add
the elapsed time $t$, discretized at resolution $\Delta$, to the state. With
$\Delta=1$\,s this multiplies the state count by at most $T_{\max}/\Delta=420$,
which gives about $3.2\times10^5$ states for a GPSR batch with the bonus bit.
\todo{Measure solve time; consider a coarser $\Delta$.}

\begin{algorithm}[t]
\caption{Interleaving by progress-vector DP}\label{alg:dp}
\begin{algorithmic}[1]
\Require segments $s_{c,j}$, order $\pre$, start pose $p_0$, gripper state $g_0$
\State $V(\mathbf{L},\ell)\gets 0$ for all $\ell$ \Comment{$\mathbf{L}=(L_1,\dots,L_n)$}
\For{$\mathbf{k}\neq\mathbf{L}$ in reverse lexicographic order}
  \State $\Lambda\gets\{0\}$ \textbf{if} $\mathbf{k}=\mathbf{0}$ \textbf{else} $\{\ell: k_\ell\ge1\}$
  \State $h\gets\textsc{Holder}(\mathbf{k},g_0)$
  \For{$\ell\in\Lambda$}
    \State $V(\mathbf{k},\ell)\gets\infty$;\; $x\gets x(\mathbf{k},\ell)$
    \For{$c$ with $k_c<L_c$}
      \State $s\gets s_{c,k_c+1}$
      \If{\textsc{PredsDone}$(s,\mathbf{k})$ \textbf{and} \textsc{GripperOK}$(s,h)$}
        \State $q\gets g(x,h,s)+V(\mathbf{k}+\mathbf{e}_c,c)$
        \If{$q<V(\mathbf{k},\ell)$}
          \State $V(\mathbf{k},\ell)\gets q$;\; $\mu(\mathbf{k},\ell)\gets c$
        \EndIf
      \EndIf
    \EndFor
  \EndFor
\EndFor
\State \Return schedule obtained by following $\mu$ from $(\mathbf{0},0)$
\end{algorithmic}
\end{algorithm}

\subsection{Execution and replanning}
The plan is executed as a behaviour tree~\todo{cite py\_trees}. Each action is
a leaf that calls the corresponding skill. Leaves report \code{RUNNING} while a
skill is active, so timeouts are checked during execution.
\todo{Implement non-blocking leaves; the current leaves block.}
Failures are handled according to the type of action:
\begin{itemize}
  \item \emph{Information actions} (\code{count}, \code{get\_visual\_info},
    \code{get\_person\_info}, \code{find\_person}) are retried once. If they
    still fail, the result register records the failure, and the robot reports
    it when the command's answer is due. The plan continues.
  \item \emph{Manipulation actions} that fail remove the remaining segments of
    their command that depend on the held object.
  \item \emph{Navigation failures} mark the location as unreachable for the
    rest of the batch.
\end{itemize}
After any change to the remaining segments, the DP is solved again from the
current pose and gripper state. Since a solve takes \X{}\,ms for a GPSR batch,
replanning does not add noticeable delay. \todo{Implement replanning; the
current system falls back to running each command from the start.}

\subsection{Optimality under failure}\label{sec:mdp}
Replanning after each failure is a receding-horizon policy that plans as if
no further failures will occur. It is not optimal for the stochastic problem.
A policy that anticipates failures may, for example, run an uncertain
perception step early to learn whether a command can be completed before
spending travel on it, avoid picking up an object just before a segment that
is likely to fail, or leave the least reliable command for last.

For GPSR batches the stochastic problem can be solved exactly. Retries within
a segment follow a fixed rule, so their effect is folded into the outcome
probabilities of the segment. Executing segment $s$ then has three outcomes:
success; an information failure, after which the segment ends and the command's
answer is marked as failed; and a manipulation failure, after which the
remaining segments of the command that depend on the object are dropped. The
segments that depend on a held object are those up to its release, which
directly follow the pick, so the dropped segments are contiguous and we model
dropping by advancing $k_c$ past them. The remaining segments of each command
are then still determined by $\mathbf{k}$. The state is $\sigma=(\mathbf{k},\ell,\mathbf{f},b)$, where
$f_c\in\{\text{ok},\text{partial},\text{failed}\}$ records the status of
command $c$ and $b$ is the interleaving bit. With $P(o\mid s,\sigma)$ the
probability of outcome $o$ and $\sigma'_o$ the resulting state, the expected
cost-to-go satisfies
\begin{equation}\label{eq:mdp}
V(\sigma)=\min_{c\in E(\sigma)}\Big[\bar g(\sigma,c)
+\sum_{o}P(o\mid s_{c,k_c+1},\sigma)\,V(\sigma'_o)\Big],
\end{equation}
where $\bar g$ is the expected per-step cost, and the terminal cost is
$-\beta\,\mathrm{Score}(\mathbf{f},b)$ with $\beta$ converting points into
seconds. Every outcome increases $\sum_c k_c$, so this Markov decision process
is acyclic, and backward induction in reverse lexicographic order of
$\mathbf{k}$ solves it exactly, as in Proposition~\ref{prop:exact}. For a GPSR
batch the state space has at most $375\cdot3^3\cdot2\approx2\times10^4$ states.
Adding elapsed time for the deadline multiplies this by $T_{\max}/\Delta$.

We use the solution of~\eqref{eq:mdp} in two ways: as an upper bound on what
any policy can achieve under the skill models of GPSR-Sim, which measures how
much replanning loses (Section~\ref{sec:gap}), and as a candidate policy for
the robot. \todo{Implement; decide whether the exact policy replaces replanning
on the robot depending on the size of the gap.}

The formulation still assumes that the locations named in a command are
correct. A person who is not where the command says, or an object that has
been moved, makes the problem partially observable, and an optimal policy
would plan over beliefs about where things are. We do not model this; the
planner acts on the locations given by the parser, and failures to find a
target are captured only through the outcome probabilities.
\todo{Note navigation failures: marking a location unreachable changes the
distance function, which the state above does not represent; these are
handled by replanning only.}

\subsection{Implementation}
The planner is a dependency-free Python module, and the executor uses
py\_trees. Both run on the robot inside a ROS~2 node. Skill implementations
are resolved by name, so the same planner and executor run unchanged in
simulation and on the robot. \todo{Robot platform description and figure.}

% ===========================================================================
\section{GPSR-Sim Benchmark}\label{sec:bench}
% ===========================================================================

\subsection{Command batches}
Each batch contains three commands sampled from the official
generator~\cite{cmdgen}. Following the rulebook, each command is rephrased by a
language model up to three times. \todo{Fix the rephrasing model and prompt.}
The generator's structured output provides the ground-truth goal for each
command. \todo{Number of batches; train/test split if any parser tuning is done.}

\subsection{Environments}
We use \X{} apartment layouts with interior walls, furniture and the standard
@Home location and object categories. \todo{Choose source: hand-built arenas,
HSSD scenes via Habitat, or both.} Distances used by the planner are path
lengths computed on the layout's occupancy grid.

\subsection{Skill models}
Each skill succeeds with probability $p_k$ and takes a duration drawn from a
distribution fitted to logs from the robot. \todo{Collect logs and fit
$p_k$ and duration distributions per skill.} People and objects can be absent
from where the command places them with probability \X.

\subsection{Evaluation}
A command counts as solved if the final world state satisfies its goal: the
object is at the requested place or with the requested person, the person has
been guided to the requested room, or the answer given to the operator matches
the ground truth. Status codes returned by skills are not used for scoring.
Each batch is scored with the GPSR score sheet, including partial credit and
the interleaving bonus.

\subsection{Fidelity tiers}
The abstract simulator runs a batch in about \X{}\,s, which allows thousands
of trials. A subset of \X{} batches is also run in \todo{Gazebo / Isaac Sim}
with the full ROS~2 navigation stack to check that conclusions carry over.

% ===========================================================================
\section{Experiments}\label{sec:exp}
% ===========================================================================

\subsection{Baselines}
\begin{itemize}
  \item \textbf{One-by-one}: the robot hears and executes each command
    separately, returning to the operator in between, as allowed by the rules.
  \item \textbf{Back-to-back}: all three commands are heard first and executed
    in the order given.
  \item \textbf{Greedy}: at each step the nearest eligible segment is chosen.
  \item \textbf{PDDL}: the parsed goals are given to Fast
    Downward~\cite{fastdownward} with a domain encoding of the skills.
    \todo{Write the domain.}
  \item \textbf{LLM merge}: a language model is given the three parsed plans
    and asked for a merged plan. \todo{Model and prompt.}
  \item \textbf{Ours (travel only)}: the DP without the dependency analysis
    or replanning, corresponding to our competition system.
  \item \textbf{Ours (full)}: dependency analysis, the cost~\eqref{eq:stage}
    and replanning after failures.
  \item \textbf{Ours (exact MDP)}: the policy from~\eqref{eq:mdp}. Under the
    simulator's skill models no policy can do better in expectation, so this
    serves as an upper bound.
\end{itemize}

\subsection{Metrics}
We report the GPSR score, the fraction of commands solved, the fraction of
batches that meet the bonus conditions, task time, distance, and the number of
constraint violations observed during execution (two objects held, an action
executed before its prerequisite, or an answer drawn from another command).

\subsection{Main results}
\begin{table*}[t]
\centering
\caption{Results on GPSR-Sim (\X{} batches, \X{} layouts). Mean $\pm$ s.d.}
\label{tab:main}
\begin{tabular}{lcccccc}
\toprule
Method & Score & Solved (\%) & Bonus (\%) & Time (s) & Distance (m) & Violations \\
\midrule
One-by-one       & \X & \X & --  & \X & \X & \X \\
Back-to-back     & \X & \X & 0   & \X & \X & \X \\
Greedy           & \X & \X & \X  & \X & \X & \X \\
PDDL             & \X & \X & \X  & \X & \X & \X \\
LLM merge        & \X & \X & \X  & \X & \X & \X \\
Ours (travel only) & \X & \X & \X & \X & \X & \X \\
Ours (full)      & \X & \X & \X  & \X & \X & \X \\
Ours (exact MDP) & \X & \X & \X  & \X & \X & \X \\
\bottomrule
\end{tabular}
\end{table*}

\todo{Describe Table~\ref{tab:main} once results are in.}

\subsection{Robustness to skill failure}
\begin{figure}[t]
\centering
\figph{4.2cm}{Placeholder: GPSR score vs.\ global skill failure rate
(0--40\%) for each method.}
\caption{Score as skill reliability decreases.}
\label{fig:failure}
\end{figure}

Figure~\ref{fig:failure} varies all success probabilities $p_k$ together.
\todo{Expected: travel-only interleaving degrades faster than one-by-one
because of restarts; replanning should close the gap. Confirm.}

\subsection{Gap to the optimal stochastic policy}\label{sec:gap}
\begin{table}[t]
\centering
\caption{Expected score of replanning relative to the exact MDP policy.}
\label{tab:gap}
\small
\begin{tabular}{lccc}
\toprule
Failure rate & Replanning & Exact MDP & Gap (\%) \\
\midrule
0\%  & \X & \X & \X \\
10\% & \X & \X & \X \\
20\% & \X & \X & \X \\
40\% & \X & \X & \X \\
\bottomrule
\end{tabular}
\end{table}
Table~\ref{tab:gap} compares replanning with the exact policy
of~\eqref{eq:mdp} at several failure rates. With no failures the two coincide,
by Proposition~\ref{prop:exact}. \todo{If the gap stays small, this justifies
replanning on the robot; if it grows with the failure rate, report which
decisions the exact policy makes differently, e.g.\ ordering risky perception
earlier.}

\subsection{Ablations}
\begin{table}[t]
\centering
\caption{Ablations of the full method.}
\label{tab:ablation}
\small
\begin{tabular}{lccc}
\toprule
Variant & Score & Time (s) & Violations \\
\midrule
Full                            & \X & \X & \X \\
-- dependency analysis          & \X & \X & \X \\
-- replanning (restart) & \X & \X & \X \\
-- gripper constraint           & \X & \X & \X \\
-- exit locations               & \X & \X & \X \\
$\kappa=0$ (no nav.\ overhead) & \X & \X & \X \\
-- bonus bit                    & \X & \X & \X \\
Euclidean distance & \X & \X & \X \\
$\lambda=0$ (time only)         & \X & \X & \X \\
\bottomrule
\end{tabular}
\end{table}

\subsection{Planner runtime}
Table~\ref{tab:runtime} compares the progress-vector DP with a subset DP
(Held-Karp over segments) on random batches of increasing size. Sizes beyond a
GPSR batch are included to show scaling.
\todo{Re-run on final implementation; current values are from a prototype.}

\begin{table}[t]
\centering
\caption{Planner runtime on random batches.}
\label{tab:runtime}
\resizebox{\columnwidth}{!}{%
\begin{tabular}{rrrrr}
\toprule
$M$ & $n$ & Progress states & Subset states ($2^M$) & Time (s) \\
\midrule
12 & 4 & 256 & 4\,096 & 0.002 \\
18 & 4 & 900 & 262\,144 & 0.007 \\
22 & 5 & 4\,500 & $4.2\times10^{6}$ & 0.022 \\
26 & 6 & 22\,500 & $6.7\times10^{7}$ & 0.30 \\
32 & 8 & 390\,625 & $4.3\times10^{9}$ & 5.9 \\
\bottomrule
\end{tabular}}
\end{table}

\subsection{Where errors come from}
\begin{table}[t]
\centering
\caption{Share of unsolved commands by first failing stage.}
\label{tab:stages}
\small
\begin{tabular}{lcc}
\toprule
Stage & Ours (travel only) & Ours (full) \\
\midrule
Speech recognition & \X & \X \\
Parsing            & \X & \X \\
Merging            & \X & \X \\
Execution          & \X & \X \\
\bottomrule
\end{tabular}
\end{table}
\todo{Annotate failures by stage; include the known case where the parser drops
the first part of a command.}

\subsection{Competition results}
\todo{RoboCup 2026 GPSR runs: commands received, plans produced, points scored,
whether the bonus was awarded, and failures observed. Include a photo or
trajectory figure.}

% ===========================================================================
\section{Discussion and Limitations}
% ===========================================================================
The dependency analysis relies on the read and write sets of each skill being
specified correctly, and on the parser naming the same object consistently
across commands. Interference that the robot cannot observe, such as a person
moving between rooms, is not modelled. The optimality results hold for the
segment decomposition and the cost and skill models; they say nothing about
errors in those models, and the exact policy of Section~\ref{sec:mdp} ignores
uncertainty about where people and objects are. The skill models in GPSR-Sim
are fitted to a single robot. \todo{Expand after results; mention whether the operator
should be consulted when commands conflict.}

% ===========================================================================
\section{Conclusion}
% ===========================================================================
We presented a planner for executing several spoken commands as one
interleaved plan, which accounts for interference between commands and
recovers from skill failures by replanning. On GPSR-Sim it \todo{one-sentence
result}. The benchmark, planner and simulator are available at \todo{URL}.

\section*{Acknowledgment}
\todo{Team, funding.}

\begin{thebibliography}{99}
\bibitem{rulebook} RoboCup@Home Technical Committee, ``RoboCup@Home rulebook,'' 2026. [Online]. Available: \url{https://github.com/RoboCupAtHome/RuleBook}
\bibitem{cmdgen} RoboCup@Home, ``GPSR command generator.'' [Online]. Available: \url{https://github.com/RoboCupAtHome/CommandGenerator}
\bibitem{saycan} M.~Ahn \emph{et al.}, ``Do as I can, not as I say: Grounding language in robotic affordances,'' arXiv:2204.01691, 2022.
\bibitem{progprompt} I.~Singh \emph{et al.}, ``ProgPrompt: Generating situated robot task plans using large language models,'' in \emph{Proc.\ ICRA}, 2023.
\bibitem{llmp} B.~Liu \emph{et al.}, ``LLM+P: Empowering large language models with optimal planning proficiency,'' arXiv:2304.11477, 2023.
\bibitem{btgenbot} R.~A.~Izzo, G.~Bardaro, and M.~Matteucci, ``BTGenBot: Behavior tree generation for robotic tasks with lightweight LLMs,'' arXiv:2403.12761, 2024.
\bibitem{frederickson1978} G.~N.~Frederickson, M.~S.~Hecht, and C.~E.~Kim, ``Approximation algorithms for some routing problems,'' \emph{SIAM J.\ Comput.}, vol.~7, no.~2, pp.~178--193, 1978.
\bibitem{savelsbergh1995} M.~W.~P.~Savelsbergh and M.~Sol, ``The general pickup and delivery problem,'' \emph{Transp.\ Sci.}, vol.~29, no.~1, pp.~17--29, 1995.
\bibitem{heldkarp1962} M.~Held and R.~M.~Karp, ``A dynamic programming approach to sequencing problems,'' \emph{J.\ SIAM}, vol.~10, no.~1, pp.~196--210, 1962.
\bibitem{nimbro2024} R.~Memmesheimer \emph{et al.}, ``RoboCup@Home 2024 OPL winner NimbRo: Anthropomorphic service robots using foundation models for perception and planning,'' arXiv:2412.14989, 2024.
\bibitem{gpsrfm2023} \todo{authors}, ``Foundation model based open vocabulary task planning and executive system for general purpose service robots,'' arXiv:2308.03357, 2023.
\bibitem{prompteng2024} \todo{authors}, ``A comparison of prompt engineering techniques for task planning and execution in service robotics,'' arXiv:2410.22997, 2024.
\bibitem{partnr} M.~Chang \emph{et al.}, ``PARTNR: A benchmark for planning and reasoning in embodied multi-agent tasks,'' arXiv:2411.00081, 2024.
\bibitem{embodiedbench} R.~Yang \emph{et al.}, ``EmbodiedBench: Comprehensive benchmarking multi-modal large language models for vision-driven embodied agents,'' arXiv:2502.09560, 2025.
\bibitem{habitat3} X.~Puig \emph{et al.}, ``Habitat 3.0: A co-habitat for humans, avatars and robots,'' arXiv:2310.13724, 2023.
\bibitem{fastdownward} M.~Helmert, ``The Fast Downward planning system,'' \emph{J.\ Artif.\ Intell.\ Res.}, vol.~26, pp.~191--246, 2006.
\end{thebibliography}

\end{document}